\subsubsection{A continuous isospectral family: measured thermal
inputs}\label{quant:heading-5}

The example extends to every unit vector \(u\in\mathbb R^3\):

\[
K_u=-(u\cdot\sigma)\otimes(u\cdot\sigma),\qquad
\mathcal H_u(G)=-\sum_{\{i,j\}}J_{ij}(u\cdot\sigma)_i(u\cdot\sigma)_j.
\]

All these Hamiltonians are related by a common single-qubit rotation and
have identical spectra on the same graph. Let \(c=u_z\), the overlap of
the interaction axis with the fixed measurement axis. If \(S\) is the
classical Ising input, the measured thermal law of \(\mathcal H_u\) is
generated by conditionally independent channels

\[
\Pr[X_i=x\mid S_i=s]=(1+cxs)/2.
\]

This follows by diagonalizing each spin along \(u\): the thermal state
is a mixture over the Ising spin configurations in that product basis,
and the squared measurement overlaps are exactly the displayed
probabilities. For normalized magnetizations,

\[
\mathbb E[M_X\mid S]=cM_S,\qquad
\operatorname{Var}(M_X\mid S)=(1-c^2)/n,
\]

so the law of total variance gives the exact identity

\[
\boxed{\operatorname{Var}(M_X)=c^2\operatorname{Var}(M_S)+(1-c^2)/n.}
\]

At fixed \(\beta>1\) on normalized complete graphs, every fixed
\(c\ne0\) therefore violates influence AA through the degree-one
magnetization, whereas \(c=0\) gives the uniform input law exactly. This
is a complete orientation classification for influence AA in this
ordered, isospectral rank-one interaction subfamily.

The positive result includes the critical point, with constants uniform
over orientations:

\[
\boxed{\max_i I_i^{\mu_{\beta,u}}(f)
\ge c_*\operatorname{Var}_{\mu_{\beta,u}}(f)^4/d^{12}
\quad(0\le\beta\le1,\ |u|=1).}
\]

Here \(c_*>0\) is absolute. To verify this, use the Curie--Weiss latent
representation proved in Section~\ref{proof:curie_weiss}. Let
\(Z\) have density proportional to
\(\exp[-n(z^2/(2\beta)-\log\cosh z)]\), and put
\(b=[n(\beta^{-1}-1)+\sqrt n]^{-1/2}\). The latent law satisfies a
Poincaré inequality with constant at most \(Cb^2\), and \(P=\tanh Z\) is
symmetric, has \(\mathbb EP^2\le Cn^{-1/2}\), and has density at least
\(c_0/b\) on \([-b/2,b/2]\).

Passing \(S\) through the measurement channel makes the measured law
exactly the product mixture with mean \(T=c\tanh Z\). If \(c=0\), it is
uniform. If \(c\ne0\), \(T\) is symmetric,
\(\mathbb ET^2\le Cn^{-1/2}\), and its density is at least
\(c_0/(|c|b)\) on \([-|c|b/2,|c|b/2]\). Thus both the low-degree moment
comparison and central-density influence comparison of Lemmas HT3--HT4
in Section~\ref{proof:high_temperature} hold with the same
absolute constants, regardless of how small \(|c|\) is.

For the spin Poincaré estimate put \(g(z)=\mathbb E_{P_{c\tanh z}}f\).
Conditional product Poincaré and differentiation give

\[
\mathbb E_Z\operatorname{Var}_{P_{c\tanh Z}}f\le\sum_i I_i^{\mu_{\beta,u}}(f),
\] \[
g'(z)=c\,\operatorname{sech}^2z\sum_i\mathbb E_{P_{c\tanh z}}D_if,
\qquad
|g'(z)|^2\le n\sum_i\mathbb E_{P_{c\tanh z}}(D_if)^2.
\]

The latent Poincaré inequality therefore yields
\(\operatorname{Var}_{\mu_{\beta,u}}f\le C\sqrt n\sum_i I_i^{\mu_{\beta,u}}(f)\).
All three inputs to Theorem HT5 hold uniformly in \(c\), so its two-case
proof gives the displayed bound unchanged. At \(\beta=0\) the law is
uniform.

The same classification holds for heat-bath influence. Conditioning on
the underlying Ising spins shows \(|\mathbb E[X_i\mid X_{-i}]|
=|c\,\mathbb E[S_i\mid X_{-i}]|\le |c|\tanh\beta\): given the other
Ising spins, the conditional mean of \(S_i\) has absolute value at most
\(\tanh\beta\), and observing their noisy outputs cannot increase this
uniform bound. Hence \(\operatorname{sech}^2\beta\,I_i\le B_i\le I_i\).

This completes the fixed-parameter diagram for normalized complete
graphs: influence AA holds \textbf{if and only if} \(\beta\le1\) or
\(u_z=0\), for fixed \(\beta\ge0\) and fixed orientation \(u\). It does
not classify arbitrary dimension-dependent approaches to this boundary.

Every measured law in this complete-graph family is exchangeable, so all
orientations and all temperatures satisfy acceptance simulation by the
theorem used in Section~\ref{quant:sec-5}.
